\subsection{Overview}

A large and internally consistent body of evidence establishes that the way a
subsurface (tile) drainage system is designed and operated governs the volume of
water it discharges, the partitioning of that water among competing
pathways---drain flow, surface runoff, evapotranspiration, and deep
seepage---and, through those hydrologic shifts, the losses of nitrogen,
phosphorus, sediment, and salts to receiving waters, as well as crop yield and
farm profitability. Three design levers recur throughout the literature: drain
spacing, drain depth, and outlet (water-table) control. Controlled drainage
(CD), also termed drainage water management (DWM), raises the outlet elevation
to hold a higher water table during non-cropping or low-demand periods, storing
water in the field and reducing cumulative drain outflow
\parencite{Singh2020,Strock2011,Blann2009}. The dominant, well-supported
mechanism is a coupling between flow and load: because most nitrate and exported
salt leave drained fields dissolved in drain water, reductions in solute loss
under controlled or shallower drainage are of roughly the same magnitude as the
reductions in outflow, with concentration changes secondary
\parencite{Wesström2007,Frankenberger2024,Lalonde1996,Tolomio2018}.

The evidence base spans multi-year paired-field experiments, controlled tank,
lysimeter, and soil-column studies, long-term monitoring, and calibrated
simulation with DRAINMOD, HYDRUS, SWAP, FLUSH, SaltMod/SahysMod, and
SWAT-family models across humid, arid, monsoonal, coastal, and cold-climate
settings, giving both process understanding and quantified outcomes. Four
consolidated conclusions frame this synthesis. First, controlled and shallow
drainage reliably reduce drain outflow and the loads of nitrate, phosphorus, and
salt it carries, redirecting rather than eliminating the stored water. Second,
subsurface drainage is a proven technology for leaching salts and reclaiming
waterlogged and saline land, lowering the water table below the root zone and
converting unproductive land into cropland. Third, drainage and its managed
variants raise crop yields---most strongly on poorly drained soils and in wet
years---while maintaining or improving water-use efficiency. Fourth, drainage
investment delivers favourable agronomic, financial, and social returns, with
net benefits maximised when disposal costs, water-use efficiency, and
environmental externalities are incorporated into design and appraisal. Meta-
analyses anchor the headline magnitudes: across 15 studies conventional drainage
volumes exceeded controlled drainage volumes by an average of 40\%, with CD
reducing nitrogen loss by 40--50\% and surface-water phosphorus loss by 25--35\%
\parencite{Strock2011}; a meta-analysis of 61 studies (1050 data pairs) found CD
decreased drainage volume by 19.2\%, reduced NO$_3$-N concentration and load,
and slightly increased crop yield (0.11\%), a ``win--win'' outcome
\parencite{Wang2020c}; a global meta-analysis of 891 experiments confirmed
drainage significantly reduces soil salinity and raises yield
\parencite{Yang2025}; and a U.S. meta-analysis of 154 comparisons found
controlled drainage and subirrigation raised yield by 8.0\% over free drainage
\parencite{Singh2022b}.

\subsection{Hydrologic Outcomes: Outflow Reduction and the Repartitioning of Water}

\subsubsection{Reduction of subsurface drain outflow}

The most robust and frequently replicated finding is that controlled drainage
reduces cumulative and annual subsurface discharge relative to free
(conventional) drainage. In Sweden, CD produced much lower total outflow in both
measured years and lowered nitrate and phosphorus losses in proportion
\parencite{Wesström2001,Wesström2007}. In Quebec, water-table control at 0.25
and 0.50~m above the drain reduced total flow by 58.7\% and 65.3\% (1992) and
40.9\% and 95\% (1993), with nitrate-loss reductions up to 95.7\% attributed
primarily to decreased flow \parencite{Lalonde1996}. Iowa comparisons found
controlled and shallow drainage cut subsurface volume by 37\% and 46\%
(four-year averages) and by 60\% and 58\% over five years
\parencite{Helmers2012,Schott2017}. Ohio field data showed 16--24\% reductions
in annual discharge \parencite{Pease2017}, and across six northwest-Ohio sites
mean daily managed-zone reductions ranged 40--100\%, with linear mixed models
explaining 57--92\% of variance \parencite{Gunn2015}. Indiana paired-field work
found the strongest reductions in the dormant season (up to 100\%) with annual
nitrate-load reductions of 15--31\% \parencite{Adeuya2012}, and event analysis
showed CD reduced event volume and peak flow and increased time to peak, with
antecedent conditions controlling response \parencite{Lahdou2019}.

Similar reductions recur worldwide: Italian water-table management reduced total
discharge by 81\% \parencite{Tolomio2018}; Egyptian and Iranian semiarid studies
reported significantly reduced outflow \parencite{Wahba2001,Jouni2018}; and
Indian saline Vertisols showed CD cutting outflow by about 64\% (11.0 versus
30.2~cm) and mean discharge from 3.10 to 1.12~mm~day$^{-1}$
\parencite{Karegoudar2019}. Michigan simulations projected 59--72\% reductions
\parencite{BinShokrana2023}, and a Michigan saturated-buffer study attributed
most performance to the CD component, which cut annual discharge by 47.1\% and
nitrate-N load by 48.0\% \parencite{Ghane2025}. In Finnish cultivated peat, CD
reduced discharge by 862~mm (thin) and 457~mm (thick peat) with only a modest
groundwater rise \parencite{Tähtikarhu2025}, and other Finnish work reported
11--23\% reductions and 17--36~cm higher water tables \parencite{Salla2022}.
Dutch reviews reported reductions of 14\% in a lysimeter and 16--85\% across
soils, with increased retention \parencite{deWit2022}, and tank experiments
confirmed the mechanism under controlled boundaries (8\% and 6\% volume
reductions at 10 and 20~cm control, drainage duration extended 17~h)
\parencite{Zheng2025}. Collectively these results establish outflow reduction,
extended drainage duration, lower and later peaks, and increased in-field
storage as consolidated, transferable effects of outlet control.

\subsubsection{Repartitioning of water among pathways}

A consolidated finding is that controlled and shallow drainage redirect rather
than eliminate stored water, most consistently into surface runoff and
secondarily into evapotranspiration and seepage. DRAINMOD simulations for Iowa
showed that at 4.0~mm~day$^{-1}$ intensity, shallow and controlled drainage each
reduced annual subsurface drainage by about 25--27~mm (roughly 15\%), most of
which reappeared as surface runoff with little change in evapotranspiration
\parencite{Singh2007b}. Belgian, North Carolina, and Canadian studies
consistently report CD decreasing subsurface drainage and increasing surface
runoff \parencite{ElSadek2002,ElSadek2002b,DEAL1986,Skaggs1995,Golmohammadi2016,Ndulue2024};
at the watershed scale in Ontario this occurred most during March--April
snowmelt, with a slight evapotranspiration increase
\parencite{Golmohammadi2016}. Polish simulations projected an approximately
fivefold increase in surface runoff from drained arable land
\parencite{Sojka2020}, and Finnish modeling found all CD methods increased
topsoil-layer runoff while raising evapotranspiration and reducing total runoff,
with design choices expressed most clearly in runoff pathways
\parencite{Isomäki2024}. Increased evapotranspiration under CD---up to a 10\%
ET$_\mathrm{act}$ increase for corn---is a recognised secondary sink
\parencite{deWit2022,DEAL1986,Ndulue2024}.

At the field-water-balance scale, the tile drain, not the soil surface, is
usually the principal conduit in humid settings. In Indiana, surface runoff was
only 1--10\% of annual precipitation versus 26--45\% for subsurface drainage
(simulated balance 7\% runoff, 34\% drain flow, 51\% evapotranspiration), with
all ponding generated by saturation excess \parencite{Saadat2020}. Iowa pothole
modeling found tile discharge and evapotranspiration accounting for 52.4\% and
47.5\% of precipitation \parencite{Amado2018,Martin2019}, and UK field drains
were the dominant pathway transferring runoff, sediment, and most sediment to
the stream, responding in every monitored event \parencite{Deasy2009}. Relative
to undrained or surface-drained conditions, however, improved subsurface
drainage generally reduces surface runoff, peak flows, sediment, and phosphorus
even as it increases nitrate delivery
\parencite{FAUSEY1995,SKAGGS1994,THOMAS1995,Gramlich2018}. In Oregon it reduced
watershed runoff and sediment yield by about 65\% and 55\%, eliminated baseflow,
and lowered peak flow \parencite{ISTOK1983}; Albanian sites recorded
surface-runoff reductions of 30--56\% \parencite{Grazhdani1996}; Minnesota
modeling found reduced runoff, soil loss, and sediment yield, especially under
no-till \parencite{Maalim2013}; and in North Carolina improved drainage cut
surface runoff by 192~mm~yr$^{-1}$ (66\%) while raising total outflow by 10\% and
strongly attenuating peak runoff (101 to 28~mm~day$^{-1}$)
\parencite{KONYHA1992}. Sandy loams showed subsurface drainage minimizing runoff
and total phosphorus loss whereas the reverse held on clay loam, phosphorus
speciation being mainly a function of soil texture \parencite{Eastman2010}.
Mechanistically, good drainage increases infiltration, lowers solute
concentrations in surface runoff \parencite{Tong2016,Tao2017}, and by lowering
the water table limits saturation-excess runoff and surface sealing
\parencite{Augeard2005}.

\subsection{Nutrient and Sediment Loss Reductions and the Flow--Load Coupling}

Because nitrate is highly mobile and travels dissolved in drain water, the
reduction in nitrate-N load under controlled drainage tracks the reduction in
outflow closely. Swedish work made this explicit: the relative decrease in N and
P loading equalled the reduction in outflow rate, with high-risk N-loss periods
coinciding with high outflow and high soil mineral N \parencite{Wesström2007}.
Global syntheses reinforce that nutrient reduction is due mainly to reduced flow
rather than consistent concentration reductions
\parencite{Frankenberger2024,Wang2020c}. A DWM synthesis drawing on 290
plot-years and 31 sites reported a mean NO$_3$-N load reduction of
13.3~kg~ha$^{-1}$~yr$^{-1}$ (46\% removal efficiency), a mean flow reduction of
126~mm~yr$^{-1}$, and dissolved reactive and total phosphorus reductions of 34\%
and 36\%, with the paired concentration reduction not significant
\parencite{Frankenberger2024}. Reported load reductions are large and
consistent: 61\% (controlled) and 49\% (shallow) in Iowa \parencite{Schott2017};
36\% and 29\% in Iowa \parencite{Helmers2012}; about 47\% at watershed scale in
Indiana, where load fell from 236.3 to 126.4~ton, mostly in the
November--March control period \parencite{Ale2012}; and 54\% for a
saturated-buffer system \parencite{Ghane2025}. In India CD roughly halved
seasonal NO$_3$-N loss (11.2 to 5.54~kg~ha$^{-1}$) \parencite{Karegoudar2019},
Italian management lowered annual N loss by 92\% (29 to 2~kg~ha$^{-1}$) and P
loss by 65\% with cumulative nitrate loss strongly correlated with cumulative
outflow \parencite{Tolomio2018}, and Chinese and Iranian studies reported drain
nitrate-loss reductions of about 15--29\% with higher soil and crop NO$_3$-N
\parencite{Dou2021,Dou2023,Jouni2018}. A broad review placed CD among effective
mitigation options (NO$_3$ reductions of 18--75\%, alongside 12--100\% for
bioreactors and 25--78\% for wetlands) \parencite{Castellano2019}.

Denitrification is a secondary, well-documented pathway where the water table is
held high: Canadian farm-scale work found greater relative abundance of
nitrate-reducing bacteria under CD and reduced suspended solids and particulate
phosphorus by damping macropore flow \parencite{Tanga2025}. Several syntheses
caution that in poorly drained soils nitrate reduction is attributable mainly to
reduced effluent rather than increased denitrification \parencite{THOMAS1995},
and Albanian data found system type barely affected concentrations, outflow
volume being the dominant control \parencite{Grazhdani1996}. Belgian modeling
reproduced an adverse relationship whereby wider spacings and higher water tables
increase denitrification and reduce leaching but also reduce plant uptake, so
total winter NO$_3$-N losses fall under CD \parencite{ElSadek2002,ElSadek2002b}.
Phosphorus responses are generally favourable but smaller and occasionally
reversed: CD lowered P loads in Sweden and Iran
\parencite{Wesström2007,Wesström2001,Jouni2018}, yet North Carolina simulations
found a small P efflux increase even as nitrate fell by up to 34\%
\parencite{DEAL1986}, and Iranian data showed slightly higher P concentrations
under CD though lower losses overall \parencite{Jouni2018}. Timing matters:
autumn--winter is the most hazardous period, and P peaks can shift toward periods
of highest outflow under CD \parencite{Tolomio2018,Wesström2007}. A recurring
caveat is that CD may increase loss through seepage and surface runoff, partially
offsetting drain-flow reductions \parencite{Frankenberger2024}.

\subsection{Design Levers: Spacing, Depth, Outlet Control, and Enhanced Configurations}

Beyond outlet control, spacing and depth systematically shift the flow--loss
balance coherently across field and simulation studies. Increasing spacing
reduces subsurface drainage and nitrate leaching to drains but increases surface
runoff, denitrification, and nitrate in runoff while decreasing plant uptake,
defining an optimum-spacing trade-off \parencite{ElSadek2002,ElSadek2002b}.
DRAINMOD-N found increasing spacing from 20 to 40~m decreased NO$_3$-N losses by
47\% while still meeting production goals \parencite{Skaggs1995}, and the MANAGE
Drain Load database (91 publications, 1279 site-years, 1961--2012) confirmed that
wider spacing and shallower placement reduce dissolved N loss, that deeper drains
show disproportionately higher N loads at a given discharge, and that the mean
and median fractions of precipitation appearing as drainage were 25\% and 20\%,
with wet years producing greater loads \parencite{Christianson2015}. Shallower
drains move less water and nutrient offsite: reducing depth from 120 to 90~cm
gave roughly 20\% reductions in annual volume and nitrate-N in Minnesota, and
76~versus~120~cm gave about 40\% \parencite{Strock2011}; regional syntheses agree
that shallow drains contain less sediment and nutrient and narrow spacing moves
more offsite \parencite{FAUSEY1995}. Chinese salinity studies identified
site-specific optima (8~m spacing at 0.7~m depth for a greenhouse soil
\parencite{Shao2012}; 0.9~m depth at 10~m spacing for chloride removal
\parencite{Yan2013b}).

Because different spacing--depth combinations achieve similar intensity, design
can be tuned to jointly meet productivity and water-quality goals. Latvian
experience over 16--17 years on about 200 fields showed similar intensity from
different combinations, with deep systematic drainage (1.3--1.5~m) giving the
best draining and economic indices in sandy and loamy soils
\parencite{SHKINKIS1979}. Belgian optimization identified a 50~m spacing at
1.25~m depth as the profit optimum but showed widening to 100~m cut total
nitrate-N loss by 20\% (37 to 29~kg~ha$^{-1}$) at small profit cost
\parencite{ElSadek2002b}. North Carolina simulations found deeper drains
increased profit but substantially increased nitrate losses, the general
prescription being to use the minimum subsurface intensity required, fit systems
to soil, and apply CD when maximum intensity is unnecessary
\parencite{Breve1998}; irrigated-system reviews likewise conclude deep drains are
often unnecessary for salinity control \parencite{Ritzema2011}.

Shallow, multi-level, and enhanced configurations extend these principles.
Shallow drainage reproduces much of the flow- and load-reduction benefit,
reducing Iowa subsurface volume by 46--58\% and nitrate by 29--49\% generally
without reducing yield \parencite{Helmers2012,Schott2017}; a multi-state analysis
found quicker water-table drawdown, reduced yield variability, and 3.4~cm less
total flow than deep drains, though with more surface runoff and a lower
benefit--cost ratio \parencite{Ghane2021b}. In saline irrigated settings drain
depth controls exported salinity: Australian deep drains flowed continuously and
exported the largest salt loads (5867~kg~ha$^{-1}$) whereas shallow drains flowed
only after irrigation, with low salinity and the greatest salt-load reduction
\parencite{Christen2001b}. Multi-level (bilevel) systems cut the time the water
table exceeded 0.5, 0.75, and 1.0~m by 84\%, 67\%, and 41\% and produced much
lower shallow-drain salinity (4.7 versus 31.7~dS~m$^{-1}$) while removing similar
total salt \parencite{Hornbuckle2007}; analytical theory shows adding shallow
drains between deep drains raises drainage rate and allows wider deep-drain
spacing \parencite{Hornbuckle2012}, and managed deep drains occupy an
intermediate position \parencite{Ayars2006b,Christen2001b}. Engineered
enhancements alter flow further: a permeable filter increased discharge (about
1.9 times conventional) and accelerated drawdown while reducing runoff
\parencite{Tao2016,Tao2017}; combined open-ditch-plus-subsurface designs improved
waterlogging removal \parencite{Ren2021,Tao2025}; siphon-type drainage raised
flow rates 27--90\% \parencite{Ren2023}; artificial macropore tubes generated
flow where none occurred and desalinated the profile \parencite{Chen2021}; and
supplementary drainage of a clay field raised mean annual discharge 2.6-fold and
cut topsoil runoff by 75\% \parencite{Häggblom2019}. Controlled drainage with
subirrigation raised groundwater levels, supplied 170--410~mm of infiltration,
and increased Dutch yields by up to 64\%
\parencite{deWit2022,deWit2024,FAUSEY1995,Almen2021,Madramootoo2007}, while land
retirement with drains closed can eliminate drainage entirely
\parencite{Purkey2001}.

\subsection{Salt Leaching and Reclamation of Saline and Waterlogged Soils}

\subsubsection{Reclamation efficacy}

Subsurface drainage reliably lowers the water table below the root zone and
removes salt from the upper profile, converting saline and waterlogged land into
productive cropland, with topsoil salinity reductions commonly 35--80\% and
corresponding yield gains, validated at full production scale
\parencite{Yang2025,Derby2024}. In the Northern Great Plains, tile-drained
topsoil salinity declined while an undrained area did not, and soybean and corn
yields improved within three years \parencite{Derby2024}. In India, drained-block
salinity fell from 7.18 to 4.62~dS~m$^{-1}$ (35.7\% mean, 16.0--66.3\% range)
while undrained rose, and yields rose 18.75--27.62\% \parencite{Sharma2006};
monsoon leaching brought weighted-average ECe to 1.4, 2.1, and 3.4~dS~m$^{-1}$ at
25, 50, and 75~m spacing \parencite{Sharma2000}; wide-spacing reclamation lowered
topsoil salinity from 8.4 to about 2.1~dS~m$^{-1}$ and raised cropping intensity
from 143\% to 191\% \parencite{Manjunatha2004}; and Tungabhadra evaluation reduced
head- and middle-reach ECe from about 9 to 1.9--3.8~dS~m$^{-1}$
\parencite{Karegoudar2022}. Chinese reclamations are similarly large: pipe
drainage cut 0--80~cm salinity from 35 to 15~g~kg$^{-1}$ versus 25--35 for open
ditches \parencite{Wang2020}; drip-plus-pipe cut EC from 14.42--15.74 to
2.86--6.91~dS~m$^{-1}$ and raised seed cotton yield 52\% \parencite{Heng2021}; and
after four years the 0--1.0~m layer fell from 27.6 to 4.1--7.2~dS~m$^{-1}$
\parencite{Zhang2021}. Desalinization is surface-concentrated: 0--20~cm rates
exceeded 80\% while deeper layers lagged \parencite{Han2023}, and the
meta-analysis found reduction greatest at $\leq$20~cm and nonsignificant below
80~cm \parencite{Yang2025}. Comparable reductions are reported for Turkey
(topsoil about 80\% at 20~cm) \parencite{Bahçeci2009}, Tunisia (mean 23.3\%)
\parencite{HajAmor2019}, and Canada (ECe 14.08 to 9.58~dS~m$^{-1}$, barley yield
up to twentyfold) \parencite{Steppuhn2006,Steppuhn2017}. Reclamation homogenises
the field (coefficient of variation 48.13\% to 28.14\%, 7~g~kg$^{-1}$ critical
salinity) \parencite{Yu2015,Yu2025} and improves soil physical condition
(porosity, infiltration 5.3 to 13.6--16.0~mm~h$^{-1}$, lower bulk density, SARe,
ESP) \parencite{Sharma2000,Yang2025}. In Egypt drainage raised yields 13--16\%
and farm income 9--44\% \parencite{Kotb2000}, whereas where the water table stays
shallow, combined waterlogging and salinity are more damaging than waterlogging
alone and yields fail above 12~dS~m$^{-1}$ \parencite{Kahlown2002}.

\subsubsection{Design parameters governing salt removal}

Closer spacing and greater depth shorten the desalinization period and increase
salt removal, though depth benefits saturate and shallow drains lose
effectiveness. A DRAINMOD analysis found the desalinization period negatively
related to intensity: at 1.2~m depth, salinity fell 55.5\% in two years at 20~m
spacing but only 5.9\% in five years at 200~m spacing, and a slightly saline
target within a decade required spacing below 30~m \parencite{Chen2025}. The
meta-analysis generalises this with climate-specific layouts---spacing $>$30~m
and depth $>$1.2~m in humid climates, 10--20~m spacing with depth $>$1.2~m in
arid climates with shallow water tables \parencite{Yang2025}. Field and column
studies identify optima near 15--20~m
\parencite{Wang2020,Wang2025c,Han2023,Chen2022,Yang2022c}, with desalination
decaying and sometimes turning negative beyond about 30~m \parencite{Zhu2025b}.
Depth effects are pronounced but bounded: pipes shallower than 1.0~m were
essentially invalid for discharging salt \parencite{Qian2021b}, and desalination
indices correlate positively with depth
\parencite{Qian2021b,Tian2023c,Tian2024}, yet SaltMod/SahysMod show diminishing
returns beyond about 1.4~m \parencite{Bahçeci2007,Bahçeci2006}. For yield under
saline irrigation an intermediate 0.8--1.2~m depth is often preferred, saving
10--15\% irrigation water with little yield penalty
\parencite{Feng2017,Feng2019,Muhammad2021}. Envelope and filter design is a
first-order control: the envelope factor dominated over pipe material (single
spun-bonded geotextile, 3.0\% perforation optimal) \parencite{Li2022}, a
double-layer geotextile integrated filtration and drainage best ($>$95\% salt
removal) while triple layers clogged \parencite{Wang2024d}, and the
O$_{90}$/d$_{90}$ ratio predicts clogging \parencite{Elzoghby2021}. Complex
stratigraphy modifies the optimum---pipes should be placed below
low-permeability silty interlayers \parencite{Guo2024b,Tian2024}---and deeper,
wider drains desalinize the whole profile whereas shallow, close drains act
mainly in the top 80~cm \parencite{KAMRA1991}. Economic analyses converge on
roughly 20~m spacing at 1.2~m depth as the maximum-return layout
\parencite{Chen2025}.

\subsubsection{Drainage-water salinity, salt balance, and temporal decline}

Because tile drains intercept salt-laden water, effluent salinity is
characteristically higher than irrigation-water salinity and declines over years
as stored salt is depleted \parencite{Li2022,Karegoudar2022}, with ECdw/ECiw
ratios of 11--17 in Iran \parencite{Akram2013} and 8000~mg~L$^{-1}$ irrigation
water producing 14{,}000~mg~L$^{-1}$ drainage water in California
\parencite{KARAJEH1994c}. Salt-balance accounting is the standard diagnostic:
salt removed exceeded salt added early on though cumulative added later exceeded
removed \parencite{Manjunatha2004}, an oasis system exported about 94\% of
imported salt while discharging 42\% of irrigation water \parencite{Askri2022},
and a straw-drainage module removed about 1.7 times the input salt
\parencite{Lu2019}; but removal can fall short with clogging (30\% then 15\%
\parencite{BenAissa2013}; no more than 27\% \parencite{HajAmor2019}). Pipes
discharged 34--46\% of soil salt-storage change during the season
\parencite{Liu2025c} and 40\% of total salt over four years \parencite{Heng2021}.
The temporal decline is a robust, long-validated pattern: most residual salt is
removed in under ten years, with effluent-salinity--drain-age relationships stable
across decades \parencite{JOHNSTON1993}, though modeling indicates 15--50 years
to full equilibrium, with drainage-water salinity more sensitive to leaching
fraction than to spacing (raising it from 0.2 to 0.4 halved equilibration time)
\parencite{Kelleners2000}. Effluent salinity typically rises then declines with a
lag of about two years \parencite{Manjunatha2004,Sharma2006,Jayawardane2001}, and
cumulative drainage governs leaching (about 70\% of salt at 1.0~m removed by
339~mm) \parencite{Bahçeci2009}. Salt export increases with both depth and
spacing (an 80~m/4~m system exported roughly 36\% more than 20~m), depth having
the larger effect where groundwater quality declines with depth
\parencite{GRISMER1993,KAMRA1991,KAMRA1994}. At watershed scale, region-wide tile
drainage raised tile salt loading from 6 to 359~kg~ha$^{-1}$ while reducing
groundwater loading, giving a modest net increase but episodic loading pulses
\parencite{Addab2022}, and the input/output salt-mass ratio stabilises near 0.4,
with higher application efficiency enhancing buildup \parencite{KARAJEH1994c}.

\subsubsection{Vertical and ionic leaching structure and hydraulics}

Leaching is vertically stratified and ionically selective. Desalinization is
strongest at the surface and weakens with depth (upper $>$ middle $>$ lower), with
deeper layers sometimes resisting or reversing
\parencite{Chen2022,Han2023,Narjary2024,Tian2024}, modulated by texture
(clayey sublayers limit leaching) \parencite{Qin2025}. Ion mobility follows
characteristic sequences (Na$^+$ $>$ Cl$^-$ $>$ SO$_4^{2-}$ $>$ Mg$^{2+}$ under
heavy rain \parencite{Yu2016}), with Cl and Na leached most efficiently
(91\% and 88\% at 0--20~cm) and Ca persisting largely as precipitated gypsum
\parencite{Qian2025,Sharma2000}. Re-accumulation is a frequent complication:
seasonal dynamics follow a ``wash-out--recovery--accumulation'' pattern
correlated with reproductive growth, some wide-spacing treatments showing
secondary-salinization risk \parencite{Zhu2025b}, late-stage re-salinization can
occur despite effective spring leaching (more severely under controlled than free
drainage) \parencite{Dou2022}, and surface accumulation reappears in dry periods
\parencite{Yu2016,Shi2021b}. Salinity around drains is spatially organised, with
desalinization zones forming within about 0.0--0.3~m of the drain and low-salinity
zones within about 5~m
\parencite{Lu2019,Heng2021,Han2023,Mirlas2003,Benyamini2005}.

A mechanistic body of work explains why leaching is often inefficient and how to
improve it. Under complete ponding, flow concentrates near the drain (about 80\%
of flow through 20\% of soil volume; full-profile desalinization required 5400~h
versus 900~h for the top 1~m) \parencite{RAO1991}, with steeply falling
strip-flow fractions \parencite{Youngs2000}. Two remedies follow: progressive or
partial flooding leaches the poorly flushed central area first (sequential
ponding used about 50\% less water \parencite{RAO1991}; a four-strip schedule used
only 1.7\% of the water and 8.3\% of the time \parencite{Youngs2000}), and
low-permeability barriers or surface mulches redistribute flow (a 10~m barrier
saved 52--82\% of leaching water \parencite{Yu2018}; a surface mulch cut leaching
time and effluent volume by 64\% and 81\% \parencite{Zhang2022}). Soil structure
governs efficiency: macropores bypass infiltration, raising time-to-leach about
twelvefold and water use eighteenfold \parencite{Xin2016}, and dual-domain
mass-transfer limitations prolong leaching \parencite{Zhang2024}. Applied-water
strategy matters: multiple low-ponding cycles matched fewer high-ponding cycles,
desalination rising with repeated cycles (about 96\% after the fourth) and
deeper ponding improving desalting \parencite{Jia2023}; more salt discharged
under flood than drip leaching, load rising linearly with discharge
\parencite{Li2022}. Ditching and deep vertical rotary tillage enhance
infiltration and desalination \parencite{Tian2023,Li2026}, whereas fine-textured
soils leach slowly and inefficiently \parencite{MANN1982}.

\subsubsection{Saline and brackish water reuse}

Subsurface drainage enables safe agricultural reuse of saline and brackish water
where salt balance is maintained by adequate leaching. Blended drainage water up
to 9~dS~m$^{-1}$ produced about 90\% wheat yield in India, cyclic use generally
outperforming continuous use \parencite{Sharma2004}; five-year data gave relative
wheat yields of 95.8, 90.3, 83.7, and 77.8\% at 6, 9, 12, and 18.8~dS~m$^{-1}$
with post-leaching root-zone salinity kept below 4~dS~m$^{-1}$
\parencite{Sharma1998}; and conjunctive-use economics favour blending
\parencite{AGNIHOTRI1992}. In China, brackish leaching discharged more salt than
imported in favourable treatments \parencite{Li2023}, saline (6~g~L$^{-1}$)
irrigation matched freshwater wheat yields with continued salinity decline
\parencite{Gao2023}, and drainage kept saline irrigation viable up to
6.25~dS~m$^{-1}$ for maize \parencite{Feng2017}, 4.0--4.4~dS~m$^{-1}$ for grain
\parencite{Feng2021}, and 8.0~dS~m$^{-1}$ under cut-soiler drainage
\parencite{Neha2022}, with cotton optima at 2--4~g~L$^{-1}$ and the preferred
depth increasing with salinity \parencite{Li2024,Lu2019}. Reuse nonetheless
carries sodicity and accumulation risks: cyclic reuse caused profile salinity
fluctuations and sodic surface soils (SARe~$>$~11)
\parencite{Mitchell2000,Jayawardane2001}, though modeled reuse over subsurface
drains caused only marginal groundwater deterioration \parencite{KARAJEH1994c}.
Blending criteria and reduced-evaporation, low-EC combinations are codified as
core salinity-control components \parencite{Kotb2000,Cuevas2019}.

Integration with complementary practices raises performance. Combined
subsurface-pipe-plus-open-ditch systems reached 170\% of a single ditch's
efficiency \parencite{Yao2024}; timing drainage to snowmelt and rain kept the
water table within 1~m and more than doubled surface salt leaching
\parencite{Yu2016}; autumn/pre-planting irrigation raised sunflower yield
9.68--14.70\% \parencite{Liu2023,Liu2025c}; windbreaks increased snow water
equivalent and, with drainage, multiplied barley yield
\parencite{Steppuhn2006,Steppuhn2017}; cut-soiler shallow drainage cut salinity
50--60\% \parencite{Neha2022,Neha2022b}; and where upward hydraulic pressure
dominates, drainage with relief wells was most effective
\parencite{Benyamini2005,Mirlas2003}. Low-cost straw-drainage modules and
cut-drains extend the approach to resource-limited settings
\parencite{Lu2019,Okuda2020}, and regional coupling of optimised drip irrigation
with tailored drainage reduced district-wide irrigation by 10\% and salinity by
15\% \parencite{Liu2026}. Italian analysis showed drainage altered redox
conditions above the drains, favouring nitrification \parencite{Mastrocicco2013};
at basin scale, however, persistent waterlogging and salinity in Central and
South Asia reflect deteriorated infrastructure and inadequate effluent disposal,
underscoring that drainage must be managed within an integrated salt balance
\parencite{Dukhovny2007,Qureshi2008,Ritzema2016,SKAGGS1994}.

\subsection{Crop Yield and Agronomic Outcomes}

\subsubsection{Magnitude of yield gains}

Direct comparisons of drained versus undrained systems repeatedly document
substantial yield increases, primarily by relieving early-season waterlogging.
Artificial drainage raised soybean yield by 275~kg~ha$^{-1}$ (8\%) in experiments
and 157~kg~ha$^{-1}$ (4\%) in producer fields, framed as effective where excess
water occurs early \parencite{Mourtzinis2021}. In a long-term Indiana experiment
corn yields progressively increased in all drained treatments across 24 corn
years while the undrained control stayed flat, though soybean was unaffected by
spacing \parencite{Rui2024,Kladivko2005}. The largest relative gains occur in
heavily stressed environments: Nile-Delta drainage doubled cotton and rice yields
and raised wheat and clover by 50\% over ten years \parencite{ELMOWELHI1988};
acid saline ``kari'' soils in India gained 1.36~t~ha$^{-1}$ (61\%)
\parencite{Mathew2001}; Bangladeshi waterlogged saline clay gained 30--64\%
(individual surface- and subsoil-drainage treatments 91\% and 81\%)
\parencite{Islam2025,Islam2022}; and an Alberta system raised barley grain yield
from 1310 to 3900~kg~ha$^{-1}$ \parencite{BENNETT1982}. More moderate but robust
gains dominate temperate systems: shallow drainage in North Carolina raised corn
by 4\% and soybean by 14.3\% \parencite{Watkins2024}; Iranian paddies exceeded
controls by 13.4--31.3\% and up to 66--70\%, the Bilevel treatment highest
\parencite{DarziNaftchali2017b,DarziNaftchali2014}; Chinese blind-ditch and pipe
drainage raised rice by 8.0\% and 6.0\% \parencite{Qi2025b}; and canola in Iran
improved roughly 60\% over surface drainage
\parencite{Asgari2021,Yang2022,Zhang2024c,SerranoCoronel2018}.

\subsubsection{Controlled drainage, subirrigation, and drainage water recycling}

Managed practices deliver benefits beyond free drainage when they retain water
for crop use in dry periods, quantified at an 8.0\% average yield increase over
free drainage across 154 comparisons \parencite{Singh2022b}. In Missouri, free
drainage and drainage water recycling raised corn yield by 20\% and 26\% over a
nondrained control (recycling gains reaching 48\% in individual years), avoiding
drought losses \parencite{Kaur2021}; in North Carolina recycling raised corn and
soybean by 8\% and 17\% with subirrigation raising the water table 15~cm
\parencite{Moursi2024}. Subirrigation performance is strongly site-dependent,
governed by lateral hydraulic conductivity: across three Ohio sites it improved
relative yield to 87--93\% and 80--91\% at two sites but had almost no effect at
Defiance owing to very low conductivity, and simulations projected it
increasingly outperforming free drainage under future climate
\parencite{Gunn2018,Gunn2019}. In Minnesota, CD and subirrigation kept soil water
optimal in both wet and dry years \parencite{Niaghi2019}; in claypan soybean
systems, drainage plus subirrigation improved yield, reduced interannual
variability, and allowed earlier planting \parencite{Nelson2011}; and combining
DWM with foliar pyraclostrobin produced synergistic soybean gains up to
1070~kg~ha$^{-1}$ \parencite{Nelson2011b}. Uzbek controlled drainage reduced
water stress between irrigations and increased grain yields
\parencite{Dukhovny2018}.

\subsubsection{Mechanisms, water-use efficiency, and irrigation savings}

The proximate mechanism is relief of excess-water stress and restoration of
root-zone aeration \parencite{BENNETT1982,DarziNaftchali2014}, with growth-stage
sensitivity (cotton flowering/boll-forming most sensitive; surface waterlogging
worse than subsurface) \parencite{Qian2021}. Improved aeration propagates through
physiology: drainage raised cotton emergence, root vigour, chlorophyll, and
dry matter, path analysis identifying emergence rate as the largest direct
contributor \parencite{Yang2022}; lowered leaf Na$^+$ and Na$^+$/K$^+$ in
sunflower \parencite{Islam2022}; raised leaf area index and growth rate in canola
\parencite{Asgari2021}; shifted rice biomass toward panicles, raising NUE by
16.1\% \parencite{DarziNaftchali2022b,Qi2025b}; and improved rhizosphere water
status through favourable root traits \parencite{Chen2025b}, while excessive
watering and no/shallow drainage increased root rot \parencite{Han2017}. Drainage
raised crop water use where growth improved (North Dakota peak ET 239 versus
183~mm; seasonal ET up 16\% and 7\%) \parencite{Rijal2012,Yang2017,Yang2022},
with an optimum water-table drawdown rate near 1.25~cm~h$^{-1}$ for cranberry
\parencite{Bulot2017}.

Controlled drainage repeatedly improves irrigation water-use efficiency without
sacrificing yield. In Egypt, simple management of existing systems improved WUE by
15--20\% (about 16\% wheat, 23\% maize, 15\% alfalfa, 20\% cotton) with no yield
loss \parencite{Wahba2005}; CD cut applied irrigation by 15.6\% (wheat) and 14\%
(maize), raised water productivity, and could save roughly 1.33~BCM for
El-Beheira \parencite{Mahmoud2021}; modified layouts reduced applied water 30--40\%
without harming salinity or yield \parencite{ELATFY1991}; and controlled-water-table
wheat achieved the highest crop water productivity at 50\% depletion with a 60~cm
water table \parencite{ElHassanin2022}. Iranian paddies raised water productivity
\parencite{DarziNaftchali2017b,Asgari2021}, Pakistani controlled drainage produced
6.5~t~ha$^{-1}$ wheat with large irrigation savings reservable for drought
\parencite{Khan2003,Sarwar2001}, and in China buried pipes raised water
productivity, highest at 1.4~m burial \parencite{Liu2025}. In arid systems the
yield benefit is often inseparable from salinity control: Nile-Delta drainage cut
top-1.5~m salinity from 5.3 to 2.2~dS~m$^{-1}$ \parencite{ELMOWELHI1988},
Kuttanad drainage leached cations and homogenised chemistry within two years
\parencite{Mathew2001}, crop-straw drains removed salt more efficiently than
plastic pipe during initial reclamation \parencite{Lu2018}, and modeling showed
supplemental irrigation $\geq$100~mm~yr$^{-1}$ kept the salt balance favourable
\parencite{Li2018,Sarwar2001,Zhang2024c,Yang2022,Liu2025,Guo2025}, though
installation without maintenance fails to control shallow water tables
\parencite{Malota2016}.

\subsubsection{Design, weather dependence, and boundary conditions}

Yield response is mediated by spacing and depth and is soil- and rainfall-
specific. Narrower spacing generally raised yield and reduced rainfall sensitivity
\parencite{Rui2024,Kladivko2005,Nelson2017}. Modeling for a Portsmouth sandy loam
required spacing of 40~m or less (30~m under poor surface drainage) for 80\%
relative yield \parencite{HARDJOAMIDJOJO1982}; across four soils, profit-maximising
spacings ranged 19--40~m near a drainage intensity of 1.1--1.2~cm~day$^{-1}$
\parencite{Skaggs2006}; Iowa DRAINMOD predicted significant gains with
soil-specific optima \parencite{Arora2024}; and Minnesota relative yield rose from
43\% undrained to 99\% at 23~m spacing and 135~cm depth \parencite{Sands2015}. In
the Canadian Prairies over-drainage was a risk---15~m spacing promoted yield loss,
40~m being a threshold avoiding negatives \parencite{Cordeiro2015}---and in North
Dakota narrower spacing suited wet years and wider spacing dry years and sugarbeet
\parencite{Acharya2019}. In China deeper burial and closer spacing generally
raised yield and desalination \parencite{Liu2025,Yang2022,Guo2025,Hou2018}, and a
0.3--1.5~m managed water table increased cane and sugar yield \parencite{CAMP1982}.

Yield benefits scale with excess-water pressure and thus with weather: only
undrained controls showed negative rainfall responses
\parencite{Rui2024,JORJANI1991}, wet years favoured narrower spacing and dry years
wider \parencite{Kladivko2005,Acharya2019,CAMP1982,Cordeiro2015}, and Japanese
yields declined with heavy-rain days only without drainage \parencite{Takeshima2023}.
Drainage also captures value through earlier field operations: it advanced sowing
about 7--8 days (worth 91--126~kg~ha$^{-1}$) \parencite{Mourtzinis2021} and avoided
29--52\% harvest population losses in wet years \parencite{Nelson2011}. The
evidence also delineates limits: where existing surface drainage is adequate,
subsurface drainage did not improve grain yield (only protein rose)
\parencite{Wiersma2010}; medium-textured soils were most responsive while fine and
coarse showed no clear effect \parencite{Singh2022b}; over-drainage can induce
drought \parencite{Cordeiro2015,DeBoer2006}; site hydraulics can nullify
subirrigation \parencite{Gunn2018,Gunn2019}; an optimum water-table depth exists
that varies with soil, crop, and permeability
\parencite{HOOGHOUDT1952,Bigelow2001}; and installation without maintenance
underperforms \parencite{Malota2016}. Drainage benefits are thus contingent on
excess-water stress being present and on appropriate, well-maintained design.
Drainage is increasingly embedded in integrated soil-and-water conservation, with
no single practice resolving difficult soils
\parencite{GARDNER1992,Manik2019,Li2021b,Rui2024}, and cover crops or perennial
forage modestly reducing early-season flow \parencite{Goeken2015,Daigh2014}, while
catchment-scale flood impacts remain less certain \parencite{Rogger2017}. Adoption
is a leading adaptive practice: about three-fourths of surveyed U.S. producers
reported drained cropland \parencite{Morton2015}, and close to 90\% of Finnish
farmers rated subsurface drainage important, ahead of irrigation
\parencite{PeltonenSainio2021}.

\subsection{Persistence of Benefits under Climate Change}

Climate-projection studies consistently conclude that the hydrologic and
water-quality benefits of controlled drainage persist, and in some seasons
strengthen, under future conditions. In Ohio, conventional discharge was
projected to decline while CD annual reductions remained near 4--6\% and summer
growing-season reductions increased, indicating greater warm-season value
\parencite{Pease2017}. Michigan simulations projected free-drainage discharge
rising 20\% by mid-century while CD continued to reduce discharge by 63--72\%,
future and historical performance being statistically indistinguishable
\parencite{BinShokrana2023}. Finnish scenarios projected CD effects remaining
similar to the historical interval \parencite{Salla2022}, and Canadian projections
of warmer, wetter conditions found CD reducing subsurface drainage more strongly
than free drainage, though rising temperature drove dry-stress yield reductions,
motivating capture, storage, and reuse of drainage water
\parencite{Ndulue2024,Gunn2018}. Chinese and Polish studies found appropriate
outlet-blocking strategies could reduce future runoff and maintain groundwater at
present levels \parencite{Awad2021,Sojka2020}, Norwegian monitoring demonstrated
rapid system response to precipitation-induced groundwater rise
\parencite{Deelstra2015}, and Polish work showed the greatest flood risk under
rapid snowmelt on near-saturated soil \parencite{Szejba2020}.

\subsection{Socio-Economic Returns and Investment Viability}

\subsubsection{Yield, income, and productivity gains}

Investment in subsurface drainage generates substantial, multi-dimensional
socio-economic returns---higher yields, cropping intensity, farm and household
income, land value, employment, and reduced out-migration
\parencite{Gopalakrishnan2007,AbdelDayem2007}. Drainage typically raises average
yields by roughly 10--30\% and lifts net farm income where leaching water and
sustainable disposal are assured \parencite{AbdelDayem2007,Datta2000,Wichelns2002}.
In the Nile Delta, \textcite{Wichelns2002} report annual gross production values
rising 500--550 per hectare and net incomes by 375, 200, and 180 per hectare in
non-saline, saline, and rehabilitation areas, with yield improvement accruing
progressively (50\% in year one to 100\% by year four). In India, the Uttar
Pradesh project reclaimed over 175{,}000~ha of sodic land, raising farm income
over 150\%, cropping intensity from 49\% to over 200\%, and household income by
53\% \parencite{Gopalakrishnan2007}; Haryana pilots increased cropping intensity
and wheat yields by 50\% to over 200\% as fallow land was cultivated and salinity
fell from about 50 to 5~dS~m$^{-1}$ \parencite{Datta2002,Datta2000}. Long-term
drainage also improves the soil resource base, raising organic carbon,
hydraulic conductivity, and enzyme activity and improving the soil quality index
by 76--136\%, with rice--wheat yields rising 328--665\% and carbon-sequestration
co-benefits up to 4.43~Mg~C~ha$^{-1}$~yr$^{-1}$
\parencite{Mukhopadhyay2023,Mukhopadhyay2023b}. In Pakistan a project lowered the
water table 40\%, raised cropping intensity from 102\% to 157\%, and lifted living
standards 30\% \parencite{Ghumman2011}, and at aggregate scale Indian waterlogging
fell over 70\% despite less than 10\% drained area \parencite{Valipour2014}.

\subsubsection{Financial returns, program scale, and cost-effectiveness}

Formal appraisals confirm favourable financial performance. The RAJAD project
(14{,}925~ha) reported a benefit--cost ratio of 1.28--2.87, payback of
4.5--6.9~years, and an internal rate of return of 18--35\%
\parencite{Gopalakrishnan2007}; the Ismila case yielded an IRR of 8--9\% rising to
about 10--11.5\%, presented as hard evidence of feasibility for small-scale,
manually installed drainage provided water for leaching and sustainable disposal
are available \parencite{Datta2000,Datta2002}. Nile-Delta drainage gave farm-level
benefit--cost ratios of 10.3 (non-saline) and 5.5 (saline) \parencite{Wichelns2002},
though not all appraisals are high (the Pakistan FDP recovered only 1.7\% per year
\parencite{Ghumman2011}). Project-planning appraisal in Texas indicated present
values of 83.1--90.9 million dollars, a 43.4 million-dollar household-income
increase, and over 46 new jobs \parencite{ROBINSON1992}, and Serbian
MODFLOW-plus-benefit--cost analysis guided phased investment
\parencite{Bozic2014}. The scale of national programmes reflects this importance:
Egyptian drainage serves about 2.4 million hectares (main), 1.9 million
(subsurface pipe), and 310{,}000 (rehabilitated) as a national success story
\parencite{AbdelDayem2007}; India reclaimed about 1.39 million ha countrywide
\parencite{Gopalakrishnan2007}; U.S. Reclamation built over 20{,}930~km of drains
protecting about 400{,}000~ha \parencite{TeKrony2004}; and globally 12--17\% of
cultivated area is drained \parencite{Valipour2014}.

Careful design and site selection sharply reduce cost. For the Nile Delta,
\textcite{Gabr2024} developed design charts giving minimum network costs of
205--374 per hectare at 300~m lateral length (37\% below 100~m), lower than
conventional approaches. Syrian optimization found vertical drainage economically
better than horizontal for its case, favouring the largest feasible spacing
\parencite{Bakour2021}. Disposal infrastructure follows economies of scale: for
Australian evaporation basins, unit cost declined with size, leakage control and
geometry being critical, with pipe drains cheaper than open drains for smaller
basins \parencite{Singh2001}, while conversion to sprinkler irrigation and
tubewell drainage with reuse gave the lowest net present cost
\parencite{Bethune2008}. Field water management modulates volumes: reducing field
length and set time was the most effective surface-irrigation upgrade, and
economic outcomes were strongly conditioned by disposal costs
\parencite{HANSON1994}. Technology lowered active-control costs: a new
flashboard-riser cut operation time and gave an average CD implementation cost of
12.93 per year per hectare \parencite{Napierala2024}, redesigning for gravity
discharge reduced pumping dependence \parencite{Imran2021}, and modernization cut
open-drain size and O\&M costs \parencite{Ritzema2023}.

\subsubsection{Materials, modernization, capacity development, and policy}

Technological change has persistently lowered costs and raised quality. Since the
1950s--60s, corrugated plastic pipes became the standard and factory pre-wrapping
of synthetic envelopes reduced installation costs by roughly 50\%
\parencite{Yannopoulos2020}; the shift from manual to mechanized drainage required
new materials, laser grade control, and quality-control systems
\parencite{Ritzema2006,AbdelDayem2007,Ritzema2007}, and Reclamation's development
of reliable design drain spacing was its most significant contribution to the
international knowledge base \parencite{TeKrony2004}. Linking applied research and
capacity development to implementation is judged ``value for money,'' with
benefits easily outweighing costs \parencite{Ritzema2007}: over 40 years of
partnership produced some of the largest, most modern programmes and documented
concrete savings---EM38 salinity measurement, three-step performance assessment,
avoidance of unjustified interceptor drains, groundwater modeling that reduced
design discharge, and lowered conservative depths---while capacity development
built mutual trust and modernized practice
\parencite{Ritzema2007,Ritzema2008,Ritzema2006}.

Because drainage water carries disposal costs and externalities, its economics
interact with policy. In California an unregulated policy produced applied-water
and drainage volumes above efficient levels, whereas a direct drainage fee induced
efficiency (society net benefit 1132~US\$/ha), with outcomes highly sensitive to
irrigation uniformity \parencite{DINAR1989}. The San Joaquin experience underscores
that economics alone do not govern outcomes: Kesterson selenium caused severe bird
mortality, leading to the conclusion that management must be technically sound,
economically viable, and socially acceptable \parencite{Letey1999}. Regional
assessments frame investment within basin-wide accounting: in Colorado combination
alternatives deepened the water table up to 1.93~m, cut soil salinity up to
950~mg/L, reduced river salt loading roughly 40--45\%, and raised relative yield up
to 9.6 points \parencite{Burkhalter2006}; the Grasslands Bypass Project cut
selenium loads 60\% \parencite{Quinn2009}; and New Zealand mitigation could reduce
N and P loads by 16\% and 23\% \parencite{McDowell2021}. However, artificial
drainage shifts diffuse phosphorus from overland to subsurface flow, giving only a
modest net reduction and acting as a connectivity pathway for preferential-flow
phosphorus \parencite{Lucci2014,Allaire2011,Dadfar2010}. Synthesising
these strands, drainage design increasingly must reconcile production with
environmental values: drainage water quality is advanced as a necessary design
criterion, and active water-table and discharge management, improved
characterization, and decision support are key future directions, since deeper
drains and wider spacing raise drain-water salinity, flow, and salt load
\parencite{Ayars2015,Yannopoulos2020,Ritzema2023}.

\subsection{Landscape Extent, Modeling, and Monitoring}

Quantifying where and how much land is drained frames the significance of these
effects. Regional mapping in Germany estimated drained proportions of about 25\%
overall, highest on hydromorphic soils (30--40\%), transferable to national scales
\parencite{Hirt2005,Hirt2011}; U.S. drainage-extent datasets diverge at watershed
scale, so best estimates depend on local conditions \parencite{Jame2022}; and
regional requirement studies documented drainable surpluses and recommended design
rates (e.g., 2~mm~day$^{-1}$) with generally good performance
\parencite{Kumar2006,Güngör2023,Ali2001,MartínezBeltrán2007,Rimidis2003}. Drain
depth is the key determinant of lateral setback for protecting wetlands
\parencite{Tangen2018}, vulnerability indices must distinguish surface from
subsurface drainage \parencite{Baffaut2020}, small-catchment and drainage-field
runoff agree well \parencite{Veinbergs2022}, and watershed models such as SWAT
under-represent subsurface drainage without explicit adaptation
\parencite{Manguerra1998,Golmohammadi2016}. Socio-economic and landscape
considerations temper further expansion in some contexts: Finnish analyses found
replacing remaining open ditches could reduce social welfare through biodiversity
and landscape losses, so surface drainage remained socially optimal for lower land
qualities \parencite{HietalaKoivu2004,Saikkonen2014}.

The quantitative understanding above rests on a mature, validated toolkit.
Physically based models reproduce field water tables and salt dynamics well:
DRAINMOD and DRAINMOD-S captured water tables (NSE up to 0.84) and salinity (NSE
0.40--0.92) and derived spacing/depth recommendations
\parencite{Chen2025,Tian2023,Tian2023c,Tian2024,Davoodi2019,Li2024}, and modified
DRAINMOD explained 63\% of corn-yield variation, with Indiana calibration reaching
Nash--Sutcliffe efficiencies of 0.70--0.85 for drain flow
\parencite{HARDJOAMIDJOJO1982,Ale2009,Wang2006b,DeBoer2006}. HYDRUS variants and
their couplings simulated water--salt transport, drainage mass, and crop response
accurately \parencite{Li2023,Dou2022,Wang2024c,Li2026,Jia2023,Liu2022,Feng2019,Feng2021},
salt-and-water-balance models (SaltMod, SahysMod, SWMS-2D, SALTIRSOIL\_M) predicted
root-zone salinity and effluent salinity closely enough for management
\parencite{Bahçeci2006,Bahçeci2007,Yao2017,Qian2021b,Abdelkrim2021}, and analytical
solutions matched barrier/mulch and dual-domain observations
\parencite{KAMRA1991,KAMRA1994,Zhang2022,Zhang2024,Youngs2000}. A simple
Hooghoudt-based method estimated profit-optimising depth and spacing
\parencite{Skaggs2006,HOOGHOUDT1952,KARAJEH1994,KARAJEH1994b,Sands2015}, watershed-
scale SWAT-Salt linked local tile behaviour to downstream salinity
\parencite{Addab2022,Liu2026,ElAgha2020}, and long-validated empirical
effluent-quality relationships remain reliable across decades
\parencite{JOHNSTON1993,Kelleners2000}. On the monitoring side,
electromagnetic-induction and apparent-EC surveys with geostatistical methods
enable rapid, nondestructive, depth-resolved assessment
\parencite{Narjary2024,Song2023,Yu2025,Singh2020d,Mirlas2003,Gates2002}, and
emerging IoT-plus-LSTM digital-twin systems achieved high-accuracy soil-EC and
water-content forecasting (test R$^2$ up to 0.97) \parencite{Qin2025}.
Collectively, this modeling and monitoring capability is what allows the design
and salt-balance principles summarised here to be transferred with confidence to
new sites, supporting the overarching conclusion that drainage systems can be
deliberately designed and managed---through spacing, depth, and outlet
control---to satisfy agricultural production and environmental objectives
simultaneously \parencite{Skaggs1995,Breve1998,ElSadek2002b,Strock2011,Ritzema2016,Blann2009,Gramlich2018,FAUSEY1995,SKAGGS1994}.